%% =============================================================
\chapter*{Sammanfattning på Svenska \texorpdfstring{\hfill\raisebox{-2em}{\includegraphics[height=3em]{img/sverige.png}}}{}}
\addcontentsline{toc}{section}{\texorpdfstring{\protect\raisebox{-0.5ex}{\protect\makebox[3ex]{\protect\includegraphics[height=1em]{img/flag-sverige.png}}}}{} in Swedish}
\index{Summary!Sammanfattning \protect\makebox[3ex]{\protect\includegraphics[height=1em]{img/flag-sverige.png}}}

Computers have been used for a variety of applications in business, science, education, engineering and so on. They help to solve real-world problems that would otherwise be slow, impossible or extremely difficult to address without computers and software. However, sometimes they do not behave exactly as we expect them. In many cases, the consequences could be very serious, for example when errors in banking or flight control software result in unexpected behaviours. Errors in computer systems are mostly not caused by the machine itself, but typically originate from the software that controls the computer systems, so-called bugs. Bugs are quite common in complex software systems since they typically have complicated input and involve many features, which makes them difficult to design and make them perfect by human effort. Detecting and fixing software bugs are important tasks in software development process. Remaining undetected bugs in any software project may lead huge problems. They can be very hard to detect and correct, especially if they are discovered after the software has been delivered. Therefore, it is very important to allocate sufficient resources, both in terms of time and manpower, to ensure that developed software is as free of bugs as possible. 

Some bugs are less serious than others. Some types of software, e.g., in user interfaces or entertainment software, can be useable even if it contains a small number of bugs.
However, in the case of critical systems and system components such as software in libraries of programming languages, 
bugs can have far-reaching consequences, and must be avoided as much as possible.
Some of libraries provide standard data structures such as stacks, queues, containers. Such data structures provide ways of storing  
and retrieving data in a way that suits the application at hand. For example, a stack allows inserting and removing elements in a particular order. Every time an element is inserted, that element is removed
in reverse order of insertion. The simplest application of a stack is to reverse a word. You insert a given word to stack - letter by letter - and then remove letters from the stack.
By using data structures data can be easily, and efficiently exchanged; it allows portability, comprehensibility, and adaptability of information.

A data structure can be both sequential or concurrent which is tricky and difficult to get correct. Concurrent data structures can be accessed and manipulated concurrently by many parallel threads are a central component of many parallel software applications. A data structure should ideally provide a simple interface to the software that uses it. An interface provides the set of operations with specifications about their types of arguments and returned values. Data structures typically use heap-allocated memory to store their data. For example, the concurrent linked queue in java.util.concurrent uses a singly linked list to organize their data. The data structure may be quite complex like skiplists and binary trees which are used to implement sets.

%\bjcom{Please extend this short paragraph to talk about complex features of data structure implementations, you can also include concurrency. Now the paragraph is too short}

The predominant method to ensure software quality is
\emph{testing}. It is a dynamic analysis where a program is executed under specific conditions, in so-called test cases, while checking whether the result for a given input matches the expected output.
%
The test cases should be carefully designed to cover as many as possible cases of program executions.
However, it is infeasible to cover all possible executions. Therefore, it is said by Edsger W. Dijkstra that testing can be used to show the presence of bugs, but never to show their absence. 
It would be nice to have techniques for checking that all executions conform to the interface of a data structure.
A possibility is to use formal techniques which is the approach used in this thesis.
%Du har säkerligen använt en dator och varit väldigt nöjd, och detta
%för att den förenklade många långa och svåra %mödosamma
%uppgifter.
%%
%Ibland gör den dock inte som du vill: du klickar på knappen ``Gör
%det'', men den följer inte ditt kommando. Eller ännu värre, datorn
%hänger upp sig och svarar inte på något kommando alls.
%%
%Viktiga arbetstider har kanske gått till spillo.
%%
%Du tar ett djupt andetag, startar om datorn, och allting fungerar igen
%som det ska.
%%
%Felet verkar komma från \emph{programvaran} som styr datorn, inte från
%själva maskinen.
%%
%Du undrar varför detta fel inte redan korrigerats, och dessutom,
%varför det från början inte \emph{försäkrades} om att felet inte kunde
%inträffa.
%
%
%Termen \emph{buggar} används generellt för att beskriva alla typer av
%fel, vare sig felet finns i dem elektriska delarna av maskinen eller i
%programmeringen.
%%
%Med tanke på hur komplexa dagens datorer är, är det inte förvånande
%att de innehåller många buggar.
%%
%De behöver ju %faktiskt
%hantera en stor mängd parametrar och egenskaper som kan generera
%% ge plats till
%många olika möjliga beteenden.
%%
%Och vi nämner inte ens den mänskliga faktorn som inför fel under
%programmeringsfasen!
%%
%Därför finns det en trend att bygga datorer med flera mindre enheter,
%som då gör datorn enklare att hantera.
%%
%Men då ställs vi inför ett nytt problem: dessa enheter kan när som
%helst kommunicera med varandra.
%%
%På grund av denna oförutsägbarhet blir det väldigt svårt att ta hänsyn
%till alla scenarier.
%
%Företag som tillverkar programvara har inget intresse av att lämna
%buggar, eftersom ett fel någonstans kan orsaka en rad andra fel i
%efterhand. %kaskad
%%
%De inkluderar därför en viktig kvalitetskontroll under
%utvecklingsfasen: om det finns för många buggar blir det helt enkelt
%inte kostnadseffektivt eller ens möjligt att senare eliminera dem.
%%
%En del buggar är dock enklare att lösa än andra.
%%
%Det spelar faktiskt ingen större roll om man inte kan svara i
%telefonen när man får ett samtal, eller om texteditorn förlorar dem
%senaste dokumentsuppdateringarna.
%%
%Detta är väldigt tråkigt, men ingen fara på taket! Vi klarar oss och
%väntar bara på programuppdateringen som löser buggen.
%%
%Å andra sidan får inga fel inträffa för dem så-kallade %
%\emph{kritiska systemen}, vars säkerhet är viktig i allra högsta
%grad. Alla fel måste elimineras, antingen i programvaran eller i
%maskinvaran.
%%
%Det är oacceptabelt, till exempel, att en pacemaker slutar fungera när
%ägaren passerar säkerhetskontrollen på flygplatsen, att krockkudden
%inte utlöses snabbt nog vid en trafikolycka om radion är på, eller om
%två tåg kolliderar på grund av att signalsystemet inte fungerat.
%%
%Därför är det nödvändigt att utveckla metoder som upptäcker
%ett felaktigt beteende i sådana system.
%
%
%Den vanligaste metoden är \emph{testing}: man kör programmet med olika
%värden för varje variabel och kollar om funktionaliteten
%överensstämmer med det som förväntas.
%%
%Dessa testscenarier är genomtänkta för att beskriva ett maximalt antal
%beteenden.
%%
%När komplexiteten däremot ökar, kan programmet befinna sig i väldigt
%många olika tillstånd. Det börjar då bli svårt att \emph{garantera}
%att metoden täcker alla möjliga programkörningar.
%%
%Det är även rent av omöjligt om systemet innehåller en parameter vars
%värde exempelvis är en siffra.
%%
%Metoden är ett effektivt sätt att upptäcka enkla buggar snabbt, men
%subtila buggar, såsom dem som kommer från en oförutsägbar tajming,
%kvarstår obemärkta.
%%
%För dem kritiska systemen där säkerheten är ytterst viktig, är det
%otänkbart att använda en metod som eventuellt tar sig igenom alla
%tester, men fortfarande innehåller buggar.
%
%För att visa att ett program verkligen uppfyller vissa egenskaper,
%måste vi först komma överens om vad det ska överensstämma med
%---vilket formellt heter en \emph{specifikation}.
%%
%Genom att lista programkonfigurationerna som borde undvikas, eller som
%önskas, alternativt båda, så beskriver man dem två kategorierna av
%specifikationer: \emph{säkerhetsegenskaper} och
%\emph{livlighetsegenskaper} (Safety och Liveness på engelska).
%%
%Till exempel, ``pacemakern stannar aldrig'', ``krockkudden måste
%öppnas inom mindre än $x$~millisekonder'', eller ``ingen enstaka
%process blockerar övriga processer'' är säkerhetsegenskaper.
%%
%Vi måste se till att programmet aldrig befinner sig i en felaktig
%konfiguration.
%% 
%Till motsatsen, ``brevbäraren levererar till rätt mottagare'',
%``systemet gör framsteg'', eller ``servern hanterar en HTTP-förfråga''
%är livlighetsegenskaper, där man observerar rätta konfigurationer,
%förknippade med vissa sannolikheter.
%% 
%Det är upp till oss att definiera vad specifikationen ska innehålla, för
%att verifiera det som önskas.
%% 
%Om ett låsningsfritt bromssystem (ABS) % Antilock Braking System
%beräknar hur hårt bilens bromsning ska vara men gör detta för sent,
%bör systemet beträktas som felaktigt.
%% 
%Denna avhandling fokuserar på säkerhetsegenskaper.
%\begin{statement}
%  {\bf Säkerhet}: %
%  {\it Givet en specifikation, kan systemet hamna i en fel konfiguration?}
%\end{statement}
%
%%% ------------------------
%Snarare än att testa programmet i specifika scenarier eller att
%analysera källkoden, fokuserar vi i denna avhandling på \emph{formell
%  verifiering}, som är ett matematiskt ramverk, för att \emph{bevisa}
%att programmet uppfyller sina egenskaper.
%
%Dessutom vill vi bevisa det helt \emph{automatiskt}, det vill säga
%utan tillsyn av användaren.
%%
%Vi börjar genom att extrahera en modell som motsvarar det ursprungliga
%programmet. I samband med det borttar vi delarna som är irrelevanta
%för den egenskapen som verifieras.
%%
%Men vad händer när programmet manipulerar obegränsade variabler? Vi
%talar då om system med ett oändligt antal tillstånd, och vi behöver
%framställa en approximation som är i linje med själva programmet.
%%
%Målet är att utforma en metod som \emph{garanterar} att inget fel
%kvarstår, och som även ger svar inom en rimlig tid.
%
%Många system med ett oändligt antal tillstånd kan faktiskt
%karakteriseras av en familj som består av system med ett ändligt antal
%tillstånd, och en parameter (eller flera) med värde från en obegränsad
%domän.
%%
%För varje värde på parametern, innehåller systemet ett ändligt antal
%tillstånd.
%%
%Parametern kan exempelvis vara
%antalet processer som är aktiva i en viss session av ett protokoll, %
%antalet noder i ett nät, %
%eller hur komponenterna av ett program kommunicerar med varandra.
%%
%Hur som helst, system som innehåller en preliminärt okänd parameter
%bör ha ett korrekt beteende oavsett värdet på parametern. De betraktas
%därför som system med ett oändligt antal tillstånd och kallas för
%parametriserade system.
%I denna avhandling presenterar vi två metoder för att verifiera vissa
%säkerhetsegenskaper hos sådana parametriserade system.
%
%Den första metoden körs baklänges. Den startar från dem felaktiga
%tillstånden, och beräknar vilka andra tillstånd som skulle kunna leda
%till ett fel.
%%
%Med andra ord upptäcker metoden alla konfigurationer som direkt eller
%indirekt är felaktiga.
%%
%Om initialtillstånden av programmet inte tillhör dem sistnämnda, anses
%programmet vara korrekt. Vi måste såklart först se till att
%approximationen av modellen motsvarar det ursprungliga programmet.
%
%Den andra metoden startar från initialtillstånden. Den begränsar sig
%till små värden av parametern, och härleder ett tröskelvärde, efter
%vilket metoden inte behöver fortsätta: den har faktiskt all nödvändig
%information för att dra slutsatsen att det inte finns några
%felkonfigurationer för större värden på parametern över denna tröskel.
%% 
%För att förenkla, bryter metoden ner varje konfiguration i små bitar
%av en viss storlek, och rekombinerar bitarna på alla möjliga sätt,
%just för att skapa konfigurationer av större storlek (litegrann som
%Lego\-bitar). %tegelstenar
%% 
%Tanken är att samla alla dem små bitarna och se till att ingen
%rekombination matchar någon felaktig konfiguration.
%% 
%I så fall är tröskeln hittad. Annars börjar man om med lite större
%bitar.
%
%
%Slutligen är det intressant att undra vilka av villkoren som behövs
%för att metodens beräkningar inte ska fortsätta för evigt. Problemet
%klassas som oavgörbart, det vill säga att det inte finns någon
%generell metod som kan lösa \emph{alla} instanser av
%problemet. Däremot har dessa två metoder trots allt visat sig vara
%väldigt effektiva.
%%
%Där andra metoder tog timmar, kunde dessa två metoder faktiskt
%verifiera vissa program inom sekunder, vilket var målet.
%%
%I vissa fall kan man även ge garanti på att antalet beräkningar är
%begränsad.

